\documentclass[11pt,a4paper]{article}
\usepackage[T1]{fontenc}
\usepackage[utf8]{inputenc}
\usepackage{lmodern}
\usepackage{amsmath,amssymb}
\usepackage[margin=25mm]{geometry}
\usepackage{longtable,booktabs,array,calc}
\usepackage[hidelinks]{hyperref}
\hypersetup{pdftitle={An action floor on transverse phase-space area produces the Yang--Mills quantum-mechanical},pdfauthor={navstokgap research working notes}}
\newcommand{\tightlist}{\setlength{\itemsep}{0pt}\setlength{\parskip}{0pt}}
\setlength{\emergencystretch}{3em}
\setcounter{secnumdepth}{0}
\begin{document}
\section{An action floor on transverse phase-space area produces the
Yang--Mills quantum-mechanical gap, and a gap forces an action
unit}\label{an-action-floor-on-transverse-phase-space-area-produces-the-yangmills-quantum-mechanical-gap-and-a-gap-forces-an-action-unit}

\begin{quote}
\textbf{Correction (2026-10-02).} A referee report and the October 1
audit (\href{corpus-audit-2026-10-01.md}{corpus audit} §3) found an
error in Proposition 5 and two weak points in the remarks. The changes,
made in place below:

\begin{enumerate}
\def\labelenumi{\arabic{enumi}.}
\tightlist
\item
  \emph{Four dimensions.} The former Proposition 5 listed the classical
  coupling \(1/g_{\rm cl}^2\) among the constants and then dropped it
  from the \(d=4\) row, concluding that four-dimensional Yang--Mills
  carries ``no constant at all'' besides \(c\) and admits no action
  unit. With the gauge field as a connection, \(1/g_{\rm cl}^2\) has the
  dimension of action in \(d=4\), and the repository's dimensionless
  coupling is \(g^2=\hbar g_{\rm cl}^2\) (\texttt{LLM.md} §2). The
  claims ``no constant at all'' and ``no action unit in \(d=4\)'' are
  \textbf{withdrawn}. Corrected row: four dimensions carry exactly one
  classical constant besides \(c\), an action unit, and no length; the
  gap cannot enter any action monomial, and the action unit with \(c\)
  (and \(\hbar\)) yields no energy, length or time. The implication that
  fails in \(d=4\) is ``action unit \(\Rightarrow\) gap''.
\item
  \emph{Classical versus quantum.} An overall factor \(1/g_{\rm cl}^2\)
  drops out of the classical field equations in every dimension. The
  former closing paragraph of Section 3 read the contrast between
  \(d=2,3\) and \(d=4\) as one between classical constants. The contrast
  becomes physical only in the quantum weight \(e^{-S_E/\hbar}\), whose
  coefficient \(1/(\hbar g_{\rm cl}^2)\) has dimension length\(^{4-d}\);
  the lead, Section 3 and Section 5 now say so.
\item
  \emph{Counting (Remark 3).} The former count kept one valley and an
  arbitrary inner cutoff \(x_0\); its coefficient \(16\sqrt2\pi m/(3g)\)
  of \(E^{3/2}\ln E\) was two thirds of the correct one. The remark now
  computes the floor-admitted volume of both valleys up to a bounded
  remainder, \((8\sqrt2\pi m/g)E^{3/2}[\ln(E/\varepsilon)+O(1)]\), and
  compares its quotient by \(h^2\) with Simon's exact leading
  asymptotics, which it matches, constant included.
\item
  \emph{Levels (Remark 4).} The Bohr--Sommerfeld estimate describes one
  channel. Counting both valleys doubles every level, the coupling near
  the origin is ignored, and the former claim of a positive first
  spacing is restricted to distinct levels.
\item
  \emph{Citation.} The Simon page ``p.~3'' was a reprint page. The
  passages are on pp.~210 and 212 of Ann. Phys. 146 (1983) 209--220; the
  constant comes from J. Funct. Anal. 53 (1983) 84--98, pp.~85--86 (both
  passage read on 2026-10-02). Remark 1's ``Lemma 0'' of G07 now points
  to the oscillator bound in the proof of G07, Theorem 5.
\end{enumerate}
\end{quote}

In the Yang--Mills quantum-mechanical model of
\href{low-dimensional-mass-gap.md}{G07} the two directions of the
analogy between a positive action floor \(h>0\) and a mass gap can both
be stated and proved inside one solved system. Forward: classical
mechanics plus the single postulate that no transverse oscillation has
phase-space area below \(h/2\) confines the escaping valley motion to
\(|x|\le2\sqrt m\,E/(\hbar g)\) at energy \(E\) and makes the admitted
phase-space volume finite,
\((8\sqrt2\pi m/g)E^{3/2}[\ln(E/\varepsilon)+O(1)]\) with both valleys
counted. Its quotient by \(h^2\) has the leading term, constant
included, of the eigenvalue counting function whose exact asymptotics
Simon proved, for every positive floor value. Bohr--Sommerfeld
quantization reproduces the energy unit
\(\varepsilon=\hbar^{4/3}g^{2/3}m^{-2/3}\) of the exact spectrum, and
the operator inequality behind G07's Theorem 5 is this postulate made
rigorous by the uncertainty inequality. The abelian model has no
transverse oscillation, so the same floor constrains nothing and its
spectrum stays continuous: with \(h>0\) assumed, a commutative field
still has no gap. Backward: in the classical model with constants \(m\)
and \(g\), any positive finite gap \(\Delta\), whatever extra constant
produces it, supplies the action unit \(\Delta^{3/4}m^{1/2}g^{-1/2}\),
so within the model ``gap'' and ``action unit'' are dimensionally
equivalent. For pure Yang--Mills the overall coupling \(1/g_{\rm cl}^2\)
drops out of the classical field equations in every dimension, so the
comparison across dimensions concerns the quantum weight
\(e^{-S_E/\hbar}\), whose coefficient \(1/(\hbar g_{\rm cl}^2)\) has
dimension length\(^{4-d}\). The equivalence holds in \(d=2\) and
\(d=3\), where that coefficient carries a length. It fails in \(d=4\),
the action-critical dimension. There \(1/g_{\rm cl}^2\) is itself an
action and the weight carries only the pure number \(1/g^2\), with
\(g^2=\hbar g_{\rm cl}^2\). An action unit therefore exists before any
gap is given and determines no energy, and a gap needs a length that
only dimensional transmutation generates. The floor postulate is a
phase-space area statement; the necessity theorem sought in STATE item 1
is the derivation of such a floor. This is G08, exploratory; ledger row
C132 cites Propositions 4--5 (Section 5).

\subsection{1. The transverse phase-space area and the floor
postulate}\label{the-transverse-phase-space-area-and-the-floor-postulate}

Take the scalar model of G07, Theorem 5, with classical Hamiltonian
\[H_{\rm cl}=\frac{p_x^2+p_y^2}{2m}+\frac{g^2}{2}x^2y^2 .\] At fixed
\(x\neq0\) the transverse motion is the oscillator
\(E_y=p_y^2/2m+\tfrac12m\omega_x^2y^2\) with \(\omega_x=g|x|/\sqrt m\).

\textbf{Proposition 1.} The transverse orbit at fixed \(x\) and
transverse energy \(E_y\) is the ellipse
\(\{p_y^2/2m+\tfrac12m\omega_x^2y^2=E_y\}\), and its enclosed
phase-space area is
\[J_y=\oint p_y\,dy=\frac{2\pi E_y}{\omega_x}=\frac{2\pi\sqrt m\,E_y}{g|x|}.\]

\emph{Proof.} The ellipse has semi-axes \(\sqrt{2mE_y}\) in \(p_y\) and
\(\sqrt{2E_y/(m\omega_x^2)}\) in \(y\); the product times \(\pi\) gives
\(2\pi E_y/\omega_x\). \(\square\)

\(J_y\) is the adiabatic invariant of the transverse oscillation: along
the valley it is conserved to leading order when \(\omega_x\) changes
slowly on the oscillation's time scale, \(|\dot x/x|\ll\omega_x\), that
is \(|x|\gg x_{\rm ad}(E)=(2E)^{1/4}g^{-1/2}\) at energy \(E\)
(Landau--Lifshitz, \emph{Mechanics} §49; the note uses only Proposition
1 and the inequality below, which need no adiabatic limit).

\textbf{Floor postulate \(\mathsf F(h)\).} No admitted state has a
transverse oscillation of phase-space area below \(h/2=\pi\hbar\):
\(J_y\ge h/2\) whenever \(x\ne0\), and \(J_x\ge h/2\) whenever
\(y\ne0\).

The value \(h/2\) is the Bohr--Sommerfeld ground-state area
\((n+\tfrac12)h\) at \(n=0\); any positive multiple of \(h\) gives the
same conclusions with a different pure number.

\subsection{2. Forward: the floor confines the valley and produces the
gap}\label{forward-the-floor-confines-the-valley-and-produces-the-gap}

\textbf{Theorem 2.} Under \(\mathsf F(h)\), every admitted classical
state of energy \(E\) satisfies
\[|x|\le X(E)=\frac{2\sqrt m\,E}{\hbar g},\qquad |y|\le X(E),\] and the
transverse energy obeys \(E_y\ge\hbar\omega_x/2=\hbar g|x|/(2\sqrt m)\).
In particular the admitted region of the energy shell is bounded,
whereas without the floor the shell contains the unbounded valley
motions \(y=p_y=0\), \(x(t)=x_0+p_xt/m\) of every energy \(p_x^2/2m>0\).

\emph{Proof.} For \(x\ne0\), Proposition 1 and \(E_y\le E\) give
\(J_y=2\pi\sqrt m\,E_y/(g|x|)\le2\pi\sqrt m\,E/(g|x|)\); the floor
\(J_y\ge\pi\hbar\) then forces \(|x|\le2\sqrt mE/(\hbar g)\). Reading
the same inequality as a bound on \(E_y\) gives
\(E_y\ge\hbar g|x|/(2\sqrt m)\). The \(y\) bound is symmetric.
\(\square\)

\textbf{Remarks.}

\begin{enumerate}
\def\labelenumi{\arabic{enumi}.}
\item
  \emph{The floor is the quantum bound.} The inequality
  \(E_y\ge\hbar g|x|/(2\sqrt m)\) is word for word the oscillator bound
  of G07, Theorem 5 (the oscillator bound in its proof):
  \(-\frac{\hbar^2}{2m}\partial_y^2
  +\frac{g^2}{2}x^2y^2\ge\hbar g|x|/(2\sqrt m)\). The quantum theorem is
  the floor postulate enforced by the uncertainty inequality
  \(\int(\mu|\phi'|^2+\nu y^2|\phi|^2)\ge\sqrt{\mu\nu}\int|\phi|^2\),
  which is the statement that no normalizable state occupies a
  transverse phase-space area below \(h/2\). So in this model the
  implication ``\(h>0\) gives the gap'' is exact, with \(h\) entering
  only through the area floor, and the rest of quantum kinematics is not
  used.
\item
  \emph{Self-consistency of the adiabatic reading.} The cutoff \(X(E)\)
  lies in the adiabatic region:
  \(X(E)/x_{\rm ad}(E)=2^{3/4}(E/\varepsilon)^{3/4}\) with
  \(\varepsilon=\hbar^{4/3}g^{2/3}m^{-2/3}\), so for \(E\gg\varepsilon\)
  the floor acts where \(J_y\) is conserved, and near the ground state,
  \(E\sim\varepsilon\), the classical reading is only an estimate.
\item
  \emph{Counting (corrected 2026-10-02).} Without the floor the
  phase-space volume \(|\{H_{\rm cl}\le E\}|\) is infinite for every
  \(E>0\) (G07, §3.1, Mechanism; Simon, Ann. Phys. 146 (1983) 209--220,
  p.~210, passage read). With the floor it is finite, and it can be
  computed with both valleys included. Write \(R=\sqrt{2E}/g\), so that
  the potential is below \(E\) exactly where \(|xy|\le R\); the momenta
  at \((x,y)\) fill a disc of area \(2\pi m(E-\tfrac{g^2}{2}x^2y^2)\).
  By Theorem 2 every admitted state lies in the box \(|x|,|y|\le X(E)\).
  For \(X(E)^2\ge R\), that is \(E\ge\varepsilon/2\), the box carries
  the volume
  \[P(E)=\iint_{|x|,|y|\le X}2\pi m\Big(E-\frac{g^2}{2}x^2y^2\Big)_+dx\,dy
  =8\pi m\,ER\Big[\frac89+\frac23\ln\frac{X^2}{R}\Big]:\] in each
  quadrant the strip \(0\le x\le R/X\), where the \(y\)-integral of
  \(E-\tfrac{g^2}2x^2y^2\) runs to \(X\), contributes \(\tfrac89ER\),
  and the region \(R/X\le x\le X\), where it runs to \(R/x\) and equals
  \(\tfrac23ER/x\), contributes \(\tfrac23ER\ln(X^2/R)\). Inside the box
  the floor removes, at each \((x,p_x)\) with \(x\ne0\), the ellipse
  \(E_y<\hbar\omega_x/2\) of area \(\pi\hbar\) in the \((y,p_y)\) plane
  (Proposition 1), and at each \((y,p_y)\) with \(y\ne0\) the ellipse
  \(E_x<\hbar\omega_y/2\) of the same area. The admitted volume
  \(V(E)=|\{H_{\rm cl}\le E\}\cap\mathsf F(h)|\) therefore obeys
  \(0\le P-V\le2\cdot\pi\hbar\cdot2X\cdot2\sqrt{2mE}=16\sqrt2\,\pi mE^{3/2}/g\).
  With \(ER=\sqrt2E^{3/2}/g\) and
  \(X^2/R=2\sqrt2\,(E/\varepsilon)^{3/2}\),
  \[V(E)=\frac{8\sqrt2\,\pi m}{g}\,E^{3/2}\Big[\ln\frac E\varepsilon+\theta(E)\Big],
  \qquad\ln2-\frac{10}9\le\theta(E)\le\ln2+\frac89 .\] Since
  \(\sqrt R=x_{\rm ad}(E)\), the hyperbola \(|xy|=R\) crosses the
  diagonal at the adiabatic scale of Remark 2. To leading order \(V\) is
  twice the one-valley integral of \((8\pi/3)\sqrt2\,mE^{3/2}/(g|x|)\),
  the transverse ellipse area \(2\pi(E-p_x^2/2m)/\omega_x\) integrated
  over \(|p_x|\le\sqrt{2mE}\), over \(x_{\rm ad}(E)\le|x|\le X(E)\). The
  former version of this remark integrated one valley from a fixed inner
  cutoff \(x_0\) and found the coefficient \(16\sqrt2\pi m/(3g)\) of
  \(E^{3/2}\ln E\), two thirds of the correct \(8\sqrt2\pi m/g\). A
  floor \(\kappa h\) with any \(\kappa>0\) in place of \(h/2\) replaces
  \(X\) by \(X/(2\kappa)\), leaves the bound on \(P-V\) unchanged (the
  excluded ellipses grow by \(2\kappa\) as the box shrinks by
  \(1/(2\kappa)\)) and shifts both bounds on \(\theta\) by
  \(-\tfrac43\ln(2\kappa)\), so the coefficient of \(E^{3/2}\ln E\) is
  independent of the floor value.

  \emph{Comparison with the exact count.} For the operator
  \(-\Delta+|xy|^\alpha\), Simon proves the limit
  \(E^{-1-1/\alpha}(\ln E)^{-1}N(E)\to\pi^{-1}\) (J. Funct. Anal. 53

  \begin{enumerate}
  \def\labelenumii{(\arabic{enumii})}
  \setcounter{enumii}{1982}
  \tightlist
  \item
    84--98, Theorem 1.4 and the display after Theorem 1.5, pp.~85--86,
    passage read). At \(\alpha=2\) this is
    \(N(E)\simeq\pi^{-1}E^{3/2}\ln E\), the growth also stated in Ann.
    Phys. 146, p.~212, end of Section 2. The dilation \(x\mapsto\ell x\)
    with \(\ell^6=\hbar^2/(mg^2)\) carries
    \(H=-\frac{\hbar^2}{2m}\Delta+\frac{g^2}{2}x^2y^2\) to
    \(\frac\varepsilon2(-\Delta+x^2y^2)\), because both coefficients
    become \(\hbar^2/(2m\ell^2)=g^2\ell^4/2=\varepsilon/2\). Hence
    \(N_H(E)\simeq(2\sqrt2/\pi)(E/\varepsilon)^{3/2}\ln(E/\varepsilon)\),
    while the floor count gives
    \(V(E)/h^2=(2\sqrt2/\pi)(E/\varepsilon)^{3/2}[\ln(E/\varepsilon)+\theta(E)]\).
    The leading terms agree, constant included. The identification
    \(N\approx V/h^2\) is the Weyl heuristic, and the agreement is a
    comparison of two computed leading terms; Simon's proof runs through
    a Feynman--Kac lower bound and sliced-bread upper bounds. Because
    the coefficient is independent of the floor value, the agreement
    tests that a positive area floor exists and leaves the number
    \(h/2\) untested. Simon (p.~86) calls a general geometric
    interpretation of the constant \(\pi^{-1}\) an open question; no
    search was run for later literature that supplies one, and no
    novelty is claimed.
  \end{enumerate}
\item
  \emph{Semiclassical levels.} With \(J_y=(n+\tfrac12)h\) conserved, the
  valley motion has effective Hamiltonian \(p_x^2/2m+\kappa_n|x|\),
  \(\kappa_n=(n+\tfrac12)\hbar g/\sqrt m\), and Bohr--Sommerfeld
  quantization \(\oint p_x\,dx=(N+\tfrac12)h\) with
  \(\oint p_x\,dx=\tfrac{8\sqrt{2m}}{3}E^{3/2}/\kappa_n\) gives
  \[E^{\rm BS}_{n,N}=\varepsilon\Big[\frac{3\pi(2n+1)(2N+1)}{16\sqrt2}\Big]^{2/3},
  \qquad\varepsilon=\hbar^{4/3}g^{2/3}m^{-2/3}.\] The unit
  \(\varepsilon\) is the exact one of G07, Theorem 5(4), as Theorem 1 of
  G07 requires; the pure numbers are semiclassical estimates of \(e_n\)
  and are labelled as such. The estimate describes one channel, the
  \(x\)-axis with its two half-valleys. The \(y\)-axis gives the same
  levels, so counting both channels doubles every level, and the
  coupling of the channels and of different \(n\) near the origin, where
  the adiabatic reading fails (Remark 2), is ignored (corrected
  2026-10-02). The estimate therefore shows a discrete spectrum with a
  positive lowest level and positive distances between distinct levels:
  in that sense the old quantum theory of this model already has a gap.
  The splitting of the doublets and the simplicity of the ground state,
  which the Jaffe--Witten conditions need, come from G07, Theorem 5(3).
  As a consistency check, each channel and each \(n\) contribute
  \(\frac{4\sqrt2}{3\pi}(E/\varepsilon)^{3/2}/(n+\tfrac12)\) levels
  below \(E\); summed over both channels and over \(n\) up to the
  adiabatic limit \(n_*\simeq2^{-1/4}(E/\varepsilon)^{3/4}\), where the
  turning point \(E/\kappa_n\) reaches \(x_{\rm ad}(E)\), they give
  \((2\sqrt2/\pi)(E/\varepsilon)^{3/2}\ln(E/\varepsilon)\) to leading
  order, the same leading term as Remark 3.
\end{enumerate}

\textbf{Corollary 3 (the commutative model).} For the abelian matrix
model, and for \(D=1\), the potential vanishes identically: there is no
transverse oscillation, \(\mathsf F(h)\) imposes no condition, the
classical energy shells stay unbounded and the quantum spectrum is
\([0,\infty)\) (G07, Theorem 6(5)). An action floor produces a gap only
in the presence of the commutator potential, which is what supplies the
transverse ellipses whose area the floor bounds.

For the SU(2) model with \(D\ge2\) the same argument applies to each of
the two transverse directions of \(\vec x_j\) relative to \(\vec x_i\):
the floor gives \(E_\perp\ge\hbar g|\vec x_i|/\sqrt m\), the bound of
G07, Theorem 6(1), and \(|\vec x_i|\le\sqrt m\,E/(\hbar g)\) on the
energy shell.

\subsection{3. Backward: a gap forces an action unit; at the
action-critical dimension the unit is classical and fixes no
scale}\label{backward-a-gap-forces-an-action-unit-at-the-action-critical-dimension-the-unit-is-classical-and-fixes-no-scale}

\textbf{Proposition 4.} Let the classical model with fixed constants
\(m\), \(g\) (dimension vectors \(d_m=(1,0,0)\),
\(d_g=(\tfrac12,-1,-1)\)) be enlarged by any fixed constants, and
suppose the enlarged theory has a positive finite gap \(\Delta\)
depending on the fixed constants only. Then the theory has the action
unit \[A=\Delta^{3/4}\,m^{1/2}\,g^{-1/2},\] and conversely an action
unit \(A\) gives the energy unit \(A^{4/3}g^{2/3}m^{-2/3}\).

\emph{Proof.} Energy has dimension vector \((1,2,-2)\) and action
\((1,2,-1)\). The system \(a\,d_m+b\,d_g+c\,(1,2,-2)=(1,2,-1)\) has the
unique solution \(a=\tfrac12\), \(b=-\tfrac12\), \(c=\tfrac34\): the
length and time components give \(-b+2c=2\), \(-b-2c=-1\), hence
\(c=\tfrac34\), \(b=-\tfrac12\), and the mass component gives
\(a=\tfrac12\). So \(m^{1/2}g^{-1/2}\Delta^{3/4}\) has action dimension
whatever the extra constants are. The converse is G07, Theorem 3 with
\(k=4\). \(\square\)

Within the model, therefore, ``the gap is positive'' and ``an action
unit exists'' are the same dimensional statement, and G07's Theorem 5
supplies the dynamical content, \(0<\delta_1<\infty\), that turns the
unit into the number \(\Delta=\delta_1\varepsilon\).

\textbf{Proposition 5 (which field theories share this equivalence;
corrected 2026-10-02).} Write the Euclidean Yang--Mills action with the
gauge field as a connection, \(D=\partial+A\),
\([A]={\rm length}^{-1}\), \(x_0=ct\):
\[S_E=\frac1{4g_{\rm cl}^2}\int F^a_{\mu\nu}F^a_{\mu\nu}\,d^dx,\qquad
[1/g_{\rm cl}^2]={\rm action}\cdot{\rm length}^{4-d}.\] The weight
\(e^{-S_E/\hbar}\) depends on the coupling through
\(\lambda_d=\hbar g_{\rm cl}^2\), of dimension length\(^{d-4}\); in
\(d=4\) the repository writes \(g^2=\hbar g_{\rm cl}^2\), a pure number
(\texttt{LLM.md} §2; G07, Proposition 10, uses the trace normalization
of the same action). With the speed \(c\) and a gap \(\Delta\):

\begin{enumerate}
\def\labelenumi{\arabic{enumi}.}
\tightlist
\item
  In \(d=2\) and \(d=3\) the pair \((1/g_{\rm cl}^2,c)\) contains
  neither an action unit nor an energy unit. Adjoining \(\Delta\) gives
  exactly one action unit, adjoining \(\hbar\) gives exactly one energy
  unit, and the equivalence of Proposition 4 holds.
\item
  In \(d=4\), \(1/g_{\rm cl}^2\) is itself an action unit. The constants
  \((1/g_{\rm cl}^2,c,\Delta)\) are dimensionally independent and their
  only action monomial is \(1/g_{\rm cl}^2\), with \(\Delta\) at
  exponent zero. The constants \((1/g_{\rm cl}^2,c,\hbar)\) contain no
  energy, length or time unit. So ``gap \(\Rightarrow\) action unit''
  holds trivially and ``action unit \(\Rightarrow\) gap'' fails.
\end{enumerate}

\begin{longtable}[]{@{}
  >{\raggedright\arraybackslash}p{(\linewidth - 6\tabcolsep) * \real{0.2500}}
  >{\raggedright\arraybackslash}p{(\linewidth - 6\tabcolsep) * \real{0.2500}}
  >{\raggedright\arraybackslash}p{(\linewidth - 6\tabcolsep) * \real{0.2500}}
  >{\raggedright\arraybackslash}p{(\linewidth - 6\tabcolsep) * \real{0.2500}}@{}}
\toprule\noalign{}
\begin{minipage}[b]{\linewidth}\raggedright
\(d\)
\end{minipage} & \begin{minipage}[b]{\linewidth}\raggedright
classical constants (with \(c\) for Yang--Mills)
\end{minipage} & \begin{minipage}[b]{\linewidth}\raggedright
action unit with \(\Delta\) adjoined
\end{minipage} & \begin{minipage}[b]{\linewidth}\raggedright
energy unit with \(\hbar\) adjoined
\end{minipage} \\
\midrule\noalign{}
\endhead
\bottomrule\noalign{}
\endlastfoot
2 & \(1/g_{\rm cl}^2\sim(1,4,-1)\) &
\((1/g_{\rm cl}^2)^{1/3}c^{-2/3}\Delta^{2/3}\) &
\(\hbar^{3/2}g_{\rm cl}\,c=\hbar c\,\lambda_2^{1/2}\) \\
3 & \(1/g_{\rm cl}^2\sim(1,3,-1)\) &
\((1/g_{\rm cl}^2)^{1/2}c^{-1/2}\Delta^{1/2}\) &
\(\hbar^2g_{\rm cl}^2\,c=\hbar c\,\lambda_3\) \\
4 & \(1/g_{\rm cl}^2\sim(1,2,-1)\), an action & \(1/g_{\rm cl}^2\)
itself; \(\Delta\) enters with exponent \(0\) & none:
\(\hbar g_{\rm cl}^2=g^2\) is a pure number \\
mechanics, \(k=-2\) & \(m\), \(\lambda\sim(1,4,-2)\) &
\(\sqrt{m\lambda}\) already classical; \(\Delta\) adds nothing & none:
\(2m\lambda/\hbar^2\) is a pure number \\
\end{longtable}

\emph{Proof.} Dimension vectors are over (mass, length, time); action is
\((1,2,-1)\), energy \((1,2,-2)\), \(c\sim(0,1,-1)\). Solve
\(a\,d_{1/g_{\rm cl}^2}+b\,d_c+c'\,(1,2,-2)=(1,2,-1)\). For \(d=3\),
adding the length and time equations gives \(2a=1\), then
\(c'=\tfrac12\), \(b=-\tfrac12\); for \(d=2\) the sum gives \(3a=1\),
then \(c'=\tfrac23\), \(b=-\tfrac23\); for \(d=4\) it gives \(a=1\),
then the mass equation gives \(c'=0\) and the length equation \(b=0\).
The determinants of the three dimension vectors are \(-3\), \(-2\),
\(-1\) for \(d=2,3,4\), so each solution is unique. With
\(\hbar\sim(1,2,-1)\) in place of \(\Delta\) the determinants are \(-2\)
and \(-1\) for \(d=2,3\), and the displayed energy monomials have vector
\((1,2,-2)\) by direct addition. In \(d=2,3\) the pair
\((1/g_{\rm cl}^2,c)\) alone fails: the mass and length equations fix
\(a=1\) and \(b=-2\) (respectively \(b=-1\)), and the time component is
then \(1\) (respectively \(0\)), never \(-1\) or \(-2\). In \(d=4\),
\(1/g_{\rm cl}^2\) and \(\hbar\) have the same vector, so the span is
that of \((1,2,-1)\) and \((0,1,-1)\); \(a(1,2,-1)+b(0,1,-1)=(1,2,-2)\)
forces \(a=1\) (mass), \(b=0\) (length) and then \(-1=-2\) (time), and
the targets \((0,1,0)\) and \((0,0,1)\) fail at the time component in
the same way. The \(k=-2\) row is G07, Section 6, Proposition 13.
\(\square\)

In \(d=2,3\) the quantum theory in infinite volume has the constants
\((\hbar,c,\lambda_d)\), which are independent, so a gap that depends on
them alone is \(\Delta=\kappa\,\hbar c\,\lambda_d^{1/(4-d)}\) with a
pure number \(\kappa\) (G07, Theorem 1). The action unit of the table is
then \(\kappa^{2/3}\hbar\) in \(d=2\) and \(\kappa^{1/2}\hbar\) in
\(d=3\): the gap returns \(\hbar\) up to a pure number.

Two facts fix the reading of the table. First, the overall factor
\(1/g_{\rm cl}^2\) drops out of the classical field equations
\(D_\mu F_{\mu\nu}=0\) in every dimension (G07, Proposition 10), and the
dilation \(A\mapsto sA(sx)\) maps solutions to solutions in every \(d\)
while multiplying the action by \(s^{4-d}\). Classical Yang--Mills
therefore has no physical scale in any dimension, and \(1/g_{\rm cl}^2\)
only converts the geometry of a solution into action and energy. The
difference between the rows becomes physical through the quantum weight
\(e^{-S_E/\hbar}\), whose coefficient
\(1/\lambda_d=1/(\hbar g_{\rm cl}^2)\) has dimension length\(^{4-d}\).
In \(d=2,3\) the weight carries a length and, with \(c\), an energy, and
a gap is a pure number times that energy. In \(d=4\) the weight carries
the pure number \(1/g^2\) and no length. Second, \(d=4\) is
action-critical in the sense of G07, Section 6, and shares its row type
with the inverse-square potential: the classical constants own an action
unit and no energy unit. This is the dimensional reason the
four-dimensional problem is one-directional. A gap needs a length that
the constants of the weight do not contain; dimensional transmutation
generates it from the regulator (G07, Proposition 10), and the gap then
returns no action unit beyond the \(1/g_{\rm cl}^2=\hbar/g^2\) already
present. Along the continuum trajectory of asymptotic freedom
\(g^2(a)\to0\), so this bare action unit grows without bound in units of
\(\hbar\); the continuum theory keeps \(\hbar\), \(c\) and the
transmuted scale, the dimensionless coupling having been traded for that
scale. The same remark applies to Proposition 4: the trajectories of
\(H_{\rm cl}\) depend on \(g^2/m\) alone and \(m\) normalizes the
action, so the equivalence there is also a statement about the quantized
model.

\subsection{4. What this decides for the positive-action
question}\label{what-this-decides-for-the-positive-action-question}

\begin{itemize}
\tightlist
\item
  \textbf{Direction.} In the solved model the analogy is an equivalence
  at the level of units and an implication at the level of dynamics: the
  floor postulate, used only through the transverse area, produces the
  gap (Theorem 2 and Remark 1), and the gap returns the unit
  (Proposition 4). The mass-gap problem is the forward direction with
  \(\hbar\) and the gauge group given; the necessity problem of STATE
  item 1 is the derivation of the floor itself. Inside this model those
  two problems have the same dimensional content, and the necessity
  problem is exactly as hard as proving a gap from a consistency premise
  that does not name \(\hbar\).
\item
  \textbf{The second ingredient.} Corollary 3 states the commutative
  case with \(h>0\) assumed: an action floor alone gives no gap; the
  commutator potential is what creates transverse oscillations for the
  floor to act on. Any transported statement ``\(h>0\) gives a gap''
  must name the structure that plays this role.
\item
  \textbf{Areas.} The floor is a bound on a phase-space area, the
  ellipse of Proposition 1. The programme's anchor, Galileo's area
  \(vF\tau^3/(6m)\) between the inertial line and the parabola, is a
  configuration-space area converted to action by the factor \(3F/v\)
  (\href{newton-insertion-action.md}{N01} §1). The model shows what a
  successful floor would have to look like on the phase-space side: a
  lower bound on the area of a closed transverse orbit, invariant under
  the adiabatic deformation of the orbit. A constant-force motion has no
  closed transverse orbit, which is one more statement of why the
  Galileo class is similarity-covariant (G07, Proposition 14).
\item
  \textbf{Gate.} For gap-track work the supplied/generated gate of G07
  can be applied in phase-space terms: a generated gap needs a family of
  closed orbits whose area can fall below \(h\) along a non-compact
  direction.
\end{itemize}

\subsection{5. Strategic consequence}\label{strategic-consequence}

G08 closes the loop inside the Yang--Mills quantum-mechanical model:
floor gives gap, gap gives unit, and the commutative model is inert
under the floor. Three facts are retained for the programme. First, the
equivalence of gap and action unit holds when the coefficient of the
quantum weight, \(1/(\hbar g_{\rm cl}^2)\), carries a length (\(d=2,3\),
and the matrix model with its constants \(m\), \(g\)). In \(d=4\) the
classical coupling is already an action unit, the weight carries a pure
number, and a gap requires a transmuted length; the contrast between the
dimensions is a quantum one, since \(1/g_{\rm cl}^2\) drops out of the
classical field equations in every dimension. (Corrected 2026-10-02: the
earlier text said \(d=4\) carries no classical constant and no action
unit.) Second, the phase-space area of a closed transverse orbit is the
object a floor must bound. Third, the floor-admitted classical volume
divided by \(h^2\) reproduces the leading eigenvalue asymptotics of the
model, constant included, for every positive floor value (Remark 3).
STATE item 1's target stays the derivation of the floor from consistency
premises; N02 gains the phase-space form of the floor as the statement
to aim at. Ledger row C132 cites Propositions 4--5; its clause ``none in
\(d=4\)'' needs the correction of Proposition 5: in \(d=4\) the
classical coupling is itself an action unit, the gap cannot enter it,
and it yields no energy unit. Theorem 2, Propositions 1, 4, 5 and the
volume formula of Remark 3 are elementary. The Weyl comparison of Remark
3 and the Bohr--Sommerfeld levels of Remark 4 are labelled estimates.
Sources are in \href{../references/batches/B79.md}{B79}; the J. Funct.
Anal. passage of Remark 3 was read on 2026-10-02 and is cited inline.
\end{document}
